% TOP_PROBLEM: {"id": 1904, "title": "Erdős Problem #52", "category_id": 1, "category": {"id": 1, "name": "number_theory", "display_name": "Number Theory", "description": "Properties of integers, prime numbers, Diophantine equations.", "slug": "number-theory", "order_index": 1, "created_at": "2026-07-31T15:26:25.670Z"}, "status": "open"}
% FIRST_POSTED: 2026-09-25
% ENTRY_KIND: statement_only
% !TeX program = lualatex
% Definition and sources only; this file is not an LLM solution attempt.
\documentclass[11pt]{article}
\usepackage[margin=25mm]{geometry}
\usepackage{fontspec}
\setmainfont{Latin Modern Roman}
\usepackage{amsmath,amssymb,mathtools}
\usepackage{unicode-math}
\setmathfont{Latin Modern Math}
\usepackage{xurl,hyperref}
\hypersetup{colorlinks=true,urlcolor=blue}
\setcounter{secnumdepth}{0}
\setlength{\parindent}{0pt}
\setlength{\parskip}{5pt}
\setlength{\emergencystretch}{3em}
\begin{document}
\section*{\texorpdfstring{Erdős Problem \#52}{Erdős Problem \#52}}\label{1904}
\subsection{Definitions and mathematical statement}
For a finite set $A\subset{\mathbb{Z}}$, write $A+A=\{a+b:a,b\in A\}$ and $A\cdot A=\{ab:a,b\in A\}$. The conjecture is
\[
 \forall\varepsilon>0\ \exists c_\varepsilon>0\ \forall A\subset{\mathbb{Z}}\text{ finite}:\qquad
 \max(|A+A|,|A\cdot A|)\ge c_\varepsilon|A|^{2-\varepsilon}.
\]
The constant is uniform in the elements and cardinality of $A$. The assertion says that addition and multiplication cannot both produce exceptionally small image sets. It does not require both sets to have nearly quadratic size.
\par\smallskip\textit{Formulation and concept references:} \href{https://www.math.tau.ac.il/~nogaa/PDFS/sumprod2.pdf}{[S1]}.

\subsection{Short English statement}
Must a finite integer set expand almost quadratically under either addition or multiplication?
\subsection{Sources}
\begin{itemize}
\item[S1]Noga Alon, Omer Angel, Itai Benjamini, and Eyal Lubetzky, Sums and products along sparse graphs, Israel Journal of Mathematics 188 (2012), 353-384.
\url{https://www.math.tau.ac.il/~nogaa/PDFS/sumprod2.pdf}.
\textit{Formulation source.}
\end{itemize}
\end{document}
